
% --- SLIDE 2: SOURCES OF NONDETERMINISM ---
\begin{frame}{Five Places Nondeterminism Hides}

    \begin{enumerate}
        \setlength\itemsep{0.45em}
        \item \textbf{Random number generators.} Weight initialisation, shuffling, dropout, data
              augmentation, train/test splits, bootstrap sampling in random forests.
        \item \textbf{Parallelism.} Floating-point addition is not associative; sum the same numbers in a
              different order on a different thread and you get a different last bit --- which compounds.
        \item \textbf{Library and hardware versions.} A minor release changes a default; a different GPU
              picks a different convolution kernel.
        \item \textbf{The data itself.} A table that is updated in place has no version. ``The training
              set'' is not a thing unless you make it one.
        \item \textbf{Configuration.} The flag you changed by hand at 11pm and did not commit.
    \end{enumerate}

    \vspace{0.8em}
    \begin{block}{Order of attack}
        Seeds and data versioning buy you the most for the least effort. Bit-exact determinism on GPU is
        the most expensive and usually the least valuable.
    \end{block}

\end{frame}
